\documentclass[11pt,a4paper]{article}
\usepackage[a4paper,margin=2.0cm]{geometry}
\usepackage{longtable,booktabs,tabularx,array}
\usepackage{enumitem,xcolor,hyperref,fancyhdr,lastpage,titlesec,parskip}
\definecolor{apblue}{HTML}{17365D}
\definecolor{aplight}{HTML}{EAF2F8}
\hypersetup{colorlinks=true,linkcolor=apblue,urlcolor=apblue,pdftitle={AP Microeconomics Lesson Plan},pdfauthor={AP-Micro-Resources}}
\pagestyle{fancy}\setlength{\headheight}{14pt}\fancyhf{}
\lhead{AP Microeconomics Lesson Plan}\rhead{APMICRO-Teaching/Lesson-Plans}\cfoot{\thepage/\pageref{LastPage}}
\titleformat{\section}{\Large\bfseries\color{apblue}}{\thesection}{0.7em}{}
\titleformat{\subsection}{\large\bfseries\color{apblue}}{\thesubsection}{0.7em}{}
\newcolumntype{Y}{>{\raggedright\arraybackslash}X}
\newcommand{\objectivebox}[1]{\noindent\colorbox{aplight}{\parbox{0.96\textwidth}{#1}}}
\newcommand{\codepath}[1]{\texttt{\small #1}}
\begin{document}
\begin{titlepage}\centering\vspace*{2cm}
{\Huge\bfseries AP Microeconomics Lesson Plan\par}\vspace{0.5cm}{\LARGE 36-session course framework and session design\par}
\vspace{1cm}{\large 18 weeks \textperiodcentered 2 meetings per week \textperiodcentered 90 minutes per meeting\par}
\vspace{0.5cm}{\large English edition\par}\vfill
\begin{tabular}{ll}Scope: & Units 1--6\\Core outputs: & concepts, graphs, calculations, explanations, FRQ writing\\Official map: & CED objective IDs and PDF page references\\Repository: & \codepath{APMICRO-Teaching/Lesson-Plans/}\\\end{tabular}
\vfill{\large \today\par}\end{titlepage}
\tableofcontents\newpage
\section{Course design}
This course treats AP Microeconomics as model-based reasoning. Students connect definitions to graphs, calculations, and causal explanations, then transfer those skills to MCQ and FRQ formats.

The default plan is 18 weeks with two 90-minute meetings per week. Every session has a visible student product and an objective mapped to the College Board CED.

\section{Course outcomes}
By the end of the course, students should be able to:
\begin{itemize}[leftmargin=2em]
\item explain scarcity, opportunity cost, comparative advantage, and the production possibilities curve;
\item use demand, supply, elasticity, surplus, trade, taxes, subsidies, and price controls to explain market outcomes;
\item calculate production, costs, revenue, profit, and profit-maximizing output;
\item compare perfect competition, monopoly, monopolistic competition, and oligopoly;
\item analyze factor markets using marginal revenue product, factor price, and monopsony models;
\item evaluate externalities, public goods, common resources, inequality, efficiency, and government policy;
\item produce AP-style graphs, calculations, and written explanations with correct labels, units, and causal logic.
\end{itemize}
\section{Standard 90-minute lesson}
\begin{tabularx}{\textwidth}{p{1.5cm}YY}\toprule Time & Phase & Student product \\ \midrule
0--8 & Retrieval warm-up & 3--5 retrieval questions or an unlabeled core graph \\
8--20 & Context and objective & A prediction and a question to answer \\
20--42 & Mini-lesson and model & Structured notes, formula, or teacher model \\
42--62 & Guided practice & Graph, calculation, or causal explanation \\
62--78 & AP-style transfer & 4--8 MCQs or one FRQ subpart \\
78--86 & Review and error code & Concept, graph, calculation, reading, or expression error \\
86--90 & Exit ticket & One conclusion and one unresolved question \\ \bottomrule\end{tabularx}
\section{Graphing and response routines}
For every graph, students identify both axes and units, name the curve that moves, mark old and new outcomes, and state the change in price, quantity, revenue, cost, surplus, or efficiency. For every calculation, they write the relationship, substitute values with units, and state the economic meaning. For every FRQ explanation, they use a because--therefore--result chain.
\section{Session-by-session design}
\subsection{Session 1: Unit 1 -- Scarcity and resource allocation}
\objectivebox{\textbf{Official Learning Objective:} Define scarce resources and explain how economic systems allocate them.\\\textbf{Objective ID:} MKT-1.A; MKT-1.B\\\textbf{CED PDF pages:} 41–42}\par\vspace{0.4em}
\begin{tabularx}{\textwidth}{p{3.0cm}Y}
\toprule
Focus & Scarcity and resource allocation\\
Student product & Use a PPF to explain one choice\\
Homework & Unit 1 HW: basic concepts\\
\bottomrule
\end{tabularx}
\subsection{Session 2: Unit 1 -- Production possibilities curve}
\objectivebox{\textbf{Official Learning Objective:} Use a PPC to explain opportunity cost, efficiency, underutilization, and growth.\\\textbf{Objective ID:} MKT-1.C\\\textbf{CED PDF pages:} 43}\par\vspace{0.4em}
\begin{tabularx}{\textwidth}{p{3.0cm}Y}
\toprule
Focus & Production possibilities curve\\
Student product & Complete a PPF and opportunity-cost explanation\\
Homework & Trade-off calculation practice\\
\bottomrule
\end{tabularx}
\subsection{Session 3: Unit 1 -- Comparative advantage and trade}
\objectivebox{\textbf{Official Learning Objective:} Determine comparative advantage and mutually beneficial specialization and trade.\\\textbf{Objective ID:} MKT-2.A; MKT-2.B\\\textbf{CED PDF pages:} 44}\par\vspace{0.4em}
\begin{tabularx}{\textwidth}{p{3.0cm}Y}
\toprule
Focus & Comparative advantage and trade\\
Student product & Complete a specialization and exchange table\\
Homework & Comparative-advantage calculations\\
\bottomrule
\end{tabularx}
\subsection{Session 4: Unit 1 -- Cost-benefit and marginal consumer choice}
\objectivebox{\textbf{Official Learning Objective:} Calculate costs and benefits and use marginal analysis to identify an optimal choice.\\\textbf{Objective ID:} CBA-1.A; CBA-1.B; CBA-2.A; CBA-2.B\\\textbf{CED PDF pages:} 45–47}\par\vspace{0.4em}
\begin{tabularx}{\textwidth}{p{3.0cm}Y}
\toprule
Focus & Cost-benefit and marginal consumer choice\\
Student product & Use MB and MC to justify a decision\\
Homework & MCQ rationale log\\
\bottomrule
\end{tabularx}
\subsection{Session 5: Unit 2 -- Demand and supply}
\objectivebox{\textbf{Official Learning Objective:} Explain demand and supply and predict how incentives and technology shift behavior.\\\textbf{Objective ID:} MKT-3.A; MKT-3.C; MKT-3.D\\\textbf{CED PDF pages:} 53; 55}\par\vspace{0.4em}
\begin{tabularx}{\textwidth}{p{3.0cm}Y}
\toprule
Focus & Demand and supply\\
Student product & Draw the initial market equilibrium\\
Homework & Unit 2 HW: equilibrium\\
\bottomrule
\end{tabularx}
\subsection{Session 6: Unit 2 -- Price elasticity of demand and supply}
\objectivebox{\textbf{Official Learning Objective:} Calculate and interpret elasticity and predict its effect on total revenue.\\\textbf{Objective ID:} MKT-3.E\\\textbf{CED PDF pages:} 56–58}\par\vspace{0.4em}
\begin{tabularx}{\textwidth}{p{3.0cm}Y}
\toprule
Focus & Price elasticity of demand and supply\\
Student product & Calculate elasticity and classify the result\\
Homework & Elasticity practice\\
\bottomrule
\end{tabularx}
\newpage
\subsection{Session 7: Unit 2 -- Other elasticities and total revenue}
\objectivebox{\textbf{Official Learning Objective:} Use income and cross-price elasticity to interpret relationships among goods and consumers.\\\textbf{Objective ID:} MKT-3.E\\\textbf{CED PDF pages:} 59}\par\vspace{0.4em}
\begin{tabularx}{\textwidth}{p{3.0cm}Y}
\toprule
Focus & Other elasticities and total revenue\\
Student product & Use elasticity to predict total-revenue changes\\
Homework & MCQ plus calculation\\
\bottomrule
\end{tabularx}
\subsection{Session 8: Unit 2 -- Market equilibrium and surplus}
\objectivebox{\textbf{Official Learning Objective:} Identify equilibrium and calculate consumer, producer, and total surplus.\\\textbf{Objective ID:} MKT-4.A\\\textbf{CED PDF pages:} 60}\par\vspace{0.4em}
\begin{tabularx}{\textwidth}{p{3.0cm}Y}
\toprule
Focus & Market equilibrium and surplus\\
Student product & Calculate consumer, producer, and total surplus\\
Homework & Surplus practice\\
\bottomrule
\end{tabularx}
\subsection{Session 9: Unit 2 -- Market disequilibrium and equilibrium changes}
\objectivebox{\textbf{Official Learning Objective:} Explain shortages, surpluses, and new equilibrium after a demand or supply change.\\\textbf{Objective ID:} MKT-4.B\\\textbf{CED PDF pages:} 61}\par\vspace{0.4em}
\begin{tabularx}{\textwidth}{p{3.0cm}Y}
\toprule
Focus & Market disequilibrium and equilibrium changes\\
Student product & Classify shortage, surplus, and the new equilibrium\\
Homework & Shift-and-equilibrium diagrams\\
\bottomrule
\end{tabularx}
\subsection{Session 10: Unit 2 -- Government intervention}
\objectivebox{\textbf{Official Learning Objective:} Graph and calculate the effects of price controls, taxes, and subsidies.\\\textbf{Objective ID:} POL-1.A\\\textbf{CED PDF pages:} 62}\par\vspace{0.4em}
\begin{tabularx}{\textwidth}{p{3.0cm}Y}
\toprule
Focus & Government intervention\\
Student product & Graph a price control, tax, or subsidy\\
Homework & FRQ subpart\\
\bottomrule
\end{tabularx}
\subsection{Session 11: Unit 2 -- International trade and public policy}
\objectivebox{\textbf{Official Learning Objective:} Explain how trade policy changes domestic outcomes and evaluate a mixed Unit 2 scenario.\\\textbf{Objective ID:} POL-1.B\\\textbf{CED PDF pages:} 63}\par\vspace{0.4em}
\begin{tabularx}{\textwidth}{p{3.0cm}Y}
\toprule
Focus & International trade and public policy\\
Student product & Complete a timed mixed Unit 2 FRQ\\
Homework & Error review\\
\bottomrule
\end{tabularx}
\subsection{Session 12: Unit 3 -- The production function}
\objectivebox{\textbf{Official Learning Objective:} Calculate and interpret total, average, and marginal product.\\\textbf{Objective ID:} PRD-1.A\\\textbf{CED PDF pages:} 69}\par\vspace{0.4em}
\begin{tabularx}{\textwidth}{p{3.0cm}Y}
\toprule
Focus & The production function\\
Student product & Complete a production table and identify MP\\
Homework & Unit 3 HW: production\\
\bottomrule
\end{tabularx}
\newpage
\subsection{Session 13: Unit 3 -- Short-run production costs}
\objectivebox{\textbf{Official Learning Objective:} Calculate fixed, variable, total, average, and marginal costs in the short run.\\\textbf{Objective ID:} PRD-1.A\\\textbf{CED PDF pages:} 70}\par\vspace{0.4em}
\begin{tabularx}{\textwidth}{p{3.0cm}Y}
\toprule
Focus & Short-run production costs\\
Student product & Calculate and graph the cost curves\\
Homework & Cost-table calculations\\
\bottomrule
\end{tabularx}
\subsection{Session 14: Unit 3 -- Long-run costs and types of profit}
\objectivebox{\textbf{Official Learning Objective:} Interpret long-run costs and distinguish accounting from economic profit.\\\textbf{Objective ID:} PRD-1.A; CBA-2.C\\\textbf{CED PDF pages:} 71–72}\par\vspace{0.4em}
\begin{tabularx}{\textwidth}{p{3.0cm}Y}
\toprule
Focus & Long-run costs and types of profit\\
Student product & Interpret long-run costs and profit types\\
Homework & Profit comparison practice\\
\bottomrule
\end{tabularx}
\subsection{Session 15: Unit 3 -- Profit maximization}
\objectivebox{\textbf{Official Learning Objective:} Use marginal analysis to determine a firm’s profit-maximizing output.\\\textbf{Objective ID:} CBA-2.D\\\textbf{CED PDF pages:} 73}\par\vspace{0.4em}
\begin{tabularx}{\textwidth}{p{3.0cm}Y}
\toprule
Focus & Profit maximization\\
Student product & Use P = MR and MC to choose output\\
Homework & Profit-maximization questions\\
\bottomrule
\end{tabularx}
\subsection{Session 16: Unit 3 -- Produce, shut down, enter, or exit}
\objectivebox{\textbf{Official Learning Objective:} Explain short-run production and long-run entry or exit decisions.\\\textbf{Objective ID:} PRD-2.A\\\textbf{CED PDF pages:} 74}\par\vspace{0.4em}
\begin{tabularx}{\textwidth}{p{3.0cm}Y}
\toprule
Focus & Produce, shut down, enter, or exit\\
Student product & Apply shutdown and entry/exit conditions\\
Homework & Decision classification\\
\bottomrule
\end{tabularx}
\subsection{Session 17: Unit 3 -- Perfect competition and review}
\objectivebox{\textbf{Official Learning Objective:} Explain how perfect competition determines market and firm outcomes.\\\textbf{Objective ID:} PRD-3.A\\\textbf{CED PDF pages:} 75}\par\vspace{0.4em}
\begin{tabularx}{\textwidth}{p{3.0cm}Y}
\toprule
Focus & Perfect competition and review\\
Student product & Complete a cost/profit synthesis graph\\
Homework & Unit 3 quiz and error review\\
\bottomrule
\end{tabularx}
\subsection{Session 18: Unit 4 -- Introduction to imperfect competition}
\objectivebox{\textbf{Official Learning Objective:} Identify sources of market power and barriers to entry.\\\textbf{Objective ID:} PRD-3.B\\\textbf{CED PDF pages:} 83}\par\vspace{0.4em}
\begin{tabularx}{\textwidth}{p{3.0cm}Y}
\toprule
Focus & Introduction to imperfect competition\\
Student product & Identify barriers to entry and market power\\
Homework & Unit 4 HW: market structures\\
\bottomrule
\end{tabularx}
\newpage
\subsection{Session 19: Unit 4 -- Monopoly}
\objectivebox{\textbf{Official Learning Objective:} Determine monopoly output, price, profit, and efficiency using graphs.\\\textbf{Objective ID:} PRD-3.B\\\textbf{CED PDF pages:} 84}\par\vspace{0.4em}
\begin{tabularx}{\textwidth}{p{3.0cm}Y}
\toprule
Focus & Monopoly\\
Student product & Graph monopoly output, price, and profit\\
Homework & Monopoly graph practice\\
\bottomrule
\end{tabularx}
\subsection{Session 20: Unit 4 -- Price discrimination}
\objectivebox{\textbf{Official Learning Objective:} Explain how price discrimination changes prices, output, and surplus.\\\textbf{Objective ID:} PRD-3.B\\\textbf{CED PDF pages:} 85}\par\vspace{0.4em}
\begin{tabularx}{\textwidth}{p{3.0cm}Y}
\toprule
Focus & Price discrimination\\
Student product & Explain how price discrimination changes surplus\\
Homework & Case analysis and MCQ\\
\bottomrule
\end{tabularx}
\subsection{Session 21: Unit 4 -- Monopolistic competition}
\objectivebox{\textbf{Official Learning Objective:} Explain product differentiation and short-run and long-run monopolistic competition.\\\textbf{Objective ID:} PRD-3.B\\\textbf{CED PDF pages:} 86}\par\vspace{0.4em}
\begin{tabularx}{\textwidth}{p{3.0cm}Y}
\toprule
Focus & Monopolistic competition\\
Student product & Compare short-run and long-run outcomes\\
Homework & FRQ subpart\\
\bottomrule
\end{tabularx}
\subsection{Session 22: Unit 4 -- Oligopoly and game theory}
\objectivebox{\textbf{Official Learning Objective:} Use a payoff table to analyze strategic behavior and equilibrium.\\\textbf{Objective ID:} PRD-3.C\\\textbf{CED PDF pages:} 87}\par\vspace{0.4em}
\begin{tabularx}{\textwidth}{p{3.0cm}Y}
\toprule
Focus & Oligopoly and game theory\\
Student product & Complete a payoff table and equilibrium explanation\\
Homework & Unit 4 quiz\\
\bottomrule
\end{tabularx}
\subsection{Session 23: Unit 4 -- Market structure review}
\objectivebox{\textbf{Official Learning Objective:} Compare imperfect market structures by power, behavior, and efficiency.\\\textbf{Objective ID:} PRD-3.B; PRD-3.C\\\textbf{CED PDF pages:} 83–87}\par\vspace{0.4em}
\begin{tabularx}{\textwidth}{p{3.0cm}Y}
\toprule
Focus & Market structure review\\
Student product & Compare imperfect market structures\\
Homework & Mixed market-structure review\\
\bottomrule
\end{tabularx}
\subsection{Session 24: Unit 5 -- Introduction to factor markets}
\objectivebox{\textbf{Official Learning Objective:} Explain factor prices, derived demand, and hiring decisions.\\\textbf{Objective ID:} PRD-4.A\\\textbf{CED PDF pages:} 95}\par\vspace{0.4em}
\begin{tabularx}{\textwidth}{p{3.0cm}Y}
\toprule
Focus & Introduction to factor markets\\
Student product & Derive MRP from the product market\\
Homework & Unit 5 HW\\
\bottomrule
\end{tabularx}
\newpage
\subsection{Session 25: Unit 5 -- Changes in factor demand and factor supply}
\objectivebox{\textbf{Official Learning Objective:} Predict how productivity, output price, and labor-supply determinants shift factor markets.\\\textbf{Objective ID:} PRD-4.B\\\textbf{CED PDF pages:} 96}\par\vspace{0.4em}
\begin{tabularx}{\textwidth}{p{3.0cm}Y}
\toprule
Focus & Changes in factor demand and factor supply\\
Student product & Graph a factor-market shift\\
Homework & Wage and employment calculations\\
\bottomrule
\end{tabularx}
\subsection{Session 26: Unit 5 -- Competitive factor markets}
\objectivebox{\textbf{Official Learning Objective:} Use MRP and the factor price to determine profit-maximizing input use.\\\textbf{Objective ID:} PRD-4.C\\\textbf{CED PDF pages:} 97–98}\par\vspace{0.4em}
\begin{tabularx}{\textwidth}{p{3.0cm}Y}
\toprule
Focus & Competitive factor markets\\
Student product & Use MRP and factor price to choose input use\\
Homework & Factor-market FRQ\\
\bottomrule
\end{tabularx}
\subsection{Session 27: Unit 5 -- Monopsony and review}
\objectivebox{\textbf{Official Learning Objective:} Explain monopsony employment and wage outcomes and compare them with competition.\\\textbf{Objective ID:} PRD-4.D\\\textbf{CED PDF pages:} 99}\par\vspace{0.4em}
\begin{tabularx}{\textwidth}{p{3.0cm}Y}
\toprule
Focus & Monopsony and review\\
Student product & Compare monopsony and competitive outcomes\\
Homework & Unit 5 synthesis\\
\bottomrule
\end{tabularx}
\subsection{Session 28: Unit 6 -- Socially efficient and inefficient market outcomes}
\objectivebox{\textbf{Official Learning Objective:} Identify efficient quantity and explain how private incentives create inefficiency and DWL.\\\textbf{Objective ID:} POL-2.A; POL-2.B\\\textbf{CED PDF pages:} 105–106}\par\vspace{0.4em}
\begin{tabularx}{\textwidth}{p{3.0cm}Y}
\toprule
Focus & Socially efficient and inefficient market outcomes\\
Student product & Identify the efficient quantity and DWL\\
Homework & Unit 6 HW\\
\bottomrule
\end{tabularx}
\subsection{Session 29: Unit 6 -- Externalities}
\objectivebox{\textbf{Official Learning Objective:} Graph private and social costs or benefits and evaluate corrective policies.\\\textbf{Objective ID:} POL-3.A; POL-3.B\\\textbf{CED PDF pages:} 107}\par\vspace{0.4em}
\begin{tabularx}{\textwidth}{p{3.0cm}Y}
\toprule
Focus & Externalities\\
Student product & Graph private and social costs or benefits\\
Homework & Policy comparison\\
\bottomrule
\end{tabularx}
\subsection{Session 30: Unit 6 -- Public and private goods}
\objectivebox{\textbf{Official Learning Objective:} Classify goods and explain free riding and common-resource overuse.\\\textbf{Objective ID:} POL-3.C\\\textbf{CED PDF pages:} 108}\par\vspace{0.4em}
\begin{tabularx}{\textwidth}{p{3.0cm}Y}
\toprule
Focus & Public and private goods\\
Student product & Classify goods and explain free riding\\
Homework & Classification practice\\
\bottomrule
\end{tabularx}
\newpage
\subsection{Session 31: Unit 6 -- Government intervention in different market structures}
\objectivebox{\textbf{Official Learning Objective:} Calculate and evaluate policy effects in competitive and imperfect markets.\\\textbf{Objective ID:} POL-4.A\\\textbf{CED PDF pages:} 109–110}\par\vspace{0.4em}
\begin{tabularx}{\textwidth}{p{3.0cm}Y}
\toprule
Focus & Government intervention in different market structures\\
Student product & Evaluate a policy in a market-structure case\\
Homework & MCQ practice\\
\bottomrule
\end{tabularx}
\subsection{Session 32: Unit 6 -- Inequality}
\objectivebox{\textbf{Official Learning Objective:} Define inequality measures and explain sources of income and wealth inequality.\\\textbf{Objective ID:} POL-5.A; POL-5.B\\\textbf{CED PDF pages:} 111}\par\vspace{0.4em}
\begin{tabularx}{\textwidth}{p{3.0cm}Y}
\toprule
Focus & Inequality\\
Student product & Explain an inequality measure and its causes\\
Homework & FRQ subpart\\
\bottomrule
\end{tabularx}
\subsection{Session 33: Unit 6 -- Market failure FRQ synthesis}
\objectivebox{\textbf{Official Learning Objective:} Integrate graphs, calculations, and policy explanations in a market-failure FRQ.\\\textbf{Objective ID:} POL-2.A; POL-2.B; POL-3.A; POL-3.B; POL-3.C; POL-4.A; POL-5.A; POL-5.B\\\textbf{CED PDF pages:} 105–111}\par\vspace{0.4em}
\begin{tabularx}{\textwidth}{p{3.0cm}Y}
\toprule
Focus & Market failure FRQ synthesis\\
Student product & Complete a timed market-failure FRQ\\
Homework & Unit 6 quiz and error review\\
\bottomrule
\end{tabularx}
\subsection{Session 34: Review -- Core graphs}
\objectivebox{\textbf{Official Learning Objective:} Draw and interpret the major AP Microeconomics graphs with complete labels.\\\textbf{Objective ID:} MKT-1.C; MKT-4.A; PRD-1.A; PRD-2.A; PRD-3.A; PRD-3.B; PRD-4.C; PRD-4.D; POL-3.A; POL-4.A\\\textbf{CED PDF pages:} CED pp. 166–173}\par\vspace{0.4em}
\begin{tabularx}{\textwidth}{p{3.0cm}Y}
\toprule
Focus & Core graphs\\
Student product & Draw and interpret six core AP Micro graphs\\
Homework & Graph error table\\
\bottomrule
\end{tabularx}
\subsection{Session 35: Review -- Timed MCQ and error analysis}
\objectivebox{\textbf{Official Learning Objective:} Use course concepts, calculations, graphs, and interpretation under time pressure.\\\textbf{Objective ID:} All course skills\\\textbf{CED PDF pages:} CED pp. 15–19}\par\vspace{0.4em}
\begin{tabularx}{\textwidth}{p{3.0cm}Y}
\toprule
Focus & Timed MCQ and error analysis\\
Student product & Complete and classify a timed MCQ set\\
Homework & Targeted follow-up practice\\
\bottomrule
\end{tabularx}
\subsection{Session 36: Review -- Timed FRQ and exam readiness}
\objectivebox{\textbf{Official Learning Objective:} Produce a complete FRQ response with accurate graphs, calculations, and causal explanations.\\\textbf{Objective ID:} All course skills\\\textbf{CED PDF pages:} CED pp. 15–19; 166–173}\par\vspace{0.4em}
\begin{tabularx}{\textwidth}{p{3.0cm}Y}
\toprule
Focus & Timed FRQ and exam readiness\\
Student product & Complete a full or half FRQ set\\
Homework & Exam-readiness plan\\
\bottomrule
\end{tabularx}
\section{Materials and assessment loop}
Use the official CED and course documents to set the content boundary, the local FRQ archive for scored writing practice, the local Unit HW collection for repetition, and the open resources for supplementary explanations and problem sets. Restricted College Board MCQ stays in AP Classroom or Bluebook and is accessed through the official route.

The recurring assessment loop is retrieval, graph or calculation practice, one AP-style transfer item, error coding, and a delayed reattempt. The error log distinguishes concept, graph, calculation, reading, and expression errors.

\section{Source and synchronization note}
The objective IDs and CED pages come from \codepath{official-learning-objectives.csv} and \codepath{session-learning-objectives.csv}. The Markdown schedule is generated from the same session table. The source CED is stored at \codepath{APMICRO-MCQ-Practice/college-board-ap-microeconomics-course-and-exam-description.pdf}.

\end{document}
